记 $\nu(G)$ 为最大匹配大小，$\tau(G)$ 为最小点覆盖大小，
$\alpha(G)$ 为最大独立集大小。

\textbf{二分图的 K\H{o}nig 定理}
\[
  \nu(G)=\tau(G),
  \qquad
  \alpha(G)=|V|-\tau(G)=|V|-\nu(G).
\]
所以二分图中求最大独立集，可先求最大匹配，再取匹配大小的补集。

\textbf{Dilworth 定理}

有限偏序中，最大反链的大小等于覆盖全部元素所需的最少链数：
\[
  \text{最大反链大小}
  =\text{最少链覆盖数}.
\]
将偏序的每个元素拆成左右两份，对 $x<y$ 连边 $x_L\to y_R$。
在这个二分图中求最大匹配 $M$，则
\[
  \text{最少链覆盖数}
  =n-|M|,
  \qquad
  \text{最大反链大小}=n-|M|.
\]
注意：偏序需要使用传递闭包；若是 DAG 最小路径覆盖，边只需使用 DAG 边，公式仍为 $n-|M|$。

\textbf{反链覆盖}

Mirsky 定理给出对偶关系：覆盖偏序所有元素的最少反链数，等于偏序中最大链的大小。
